\documentclass[11pt,a4paper]{article}
\usepackage[utf8]{inputenc}\usepackage[T1]{fontenc}\usepackage[spanish,es-tabla]{babel}
\usepackage{geometry}\usepackage{graphicx}\usepackage{booktabs}\usepackage{float}\usepackage{xcolor}\usepackage{fancyhdr}\usepackage{array}\usepackage{amsmath}
\geometry{margin=2.2cm}\definecolor{hbaccent}{HTML}{1F9D8F}\definecolor{hbprimary}{HTML}{176B73}
\pagestyle{fancy}\fancyhf{}\lhead{HydroBasin Studio -- Estudio Hidrológico}\rhead{Informe Técnico}\cfoot{\thepage}

\renewcommand{\listfigurename}{Índice de Figuras}
\renewcommand{\listtablename}{Índice de Tablas}

\begin{document}

% ============================== PORTADA FORMAL ==============================
\begin{titlepage}
\vspace*{0.8cm}
\noindent
\begin{tabular}{|p{4.0cm}|p{8.2cm}|p{3.4cm}|}\hline
\textbf{\color{hbaccent}HYDROBASIN STUDIO} & \textbf{ESTUDIO HIDROLÓGICO Y MORFOMÉTRICO} \newline \small Zarzal & \small Versión 1.0 \newline 20/08/2026 \\\hline
\end{tabular}\\[2.0cm]

\begin{center}
{\Huge\bfseries INFORME DE DELIMITACIÓN Y CARACTERIZACIÓN MORFOMÉTRICA DE CUENCA}\\[0.6cm]
{\Large\bfseries Zarzal}\\[0.3cm]
{\normalsize Barrancabermeja, Santander, Colombia}\\[2.2cm]

\textbf{RELACIÓN DE REVISIÓN Y AUTORÍA}\\[3mm]
\begin{tabular}{|p{4.5cm}|p{5.0cm}|p{4.5cm}|}\hline
\textbf{ELABORÓ:} & \textbf{REVISÓ:} & \textbf{FECHA:} \\\hline
Jana Criollo & Mario Díaz & 20/08/2026 \\\hline
\multicolumn{2}{|l|}{\textbf{CLIENTE / ENTIDAD:} IDOM} & \textbf{ESTADO:} Aprobado \\\hline
\end{tabular}
\end{center}

\vfill
\noindent
{\footnotesize Documento técnico de ingeniería hidrológica generado por HydroBasin Studio. Prohibida su alteración no autorizada.}
\end{titlepage}

% ============================== ÍNDICES ==============================
\tableofcontents
\vspace{0.6cm}
\listoffigures
\vspace{0.6cm}
\listoftables
\newpage

\section{Introducción}
El presente documento corresponde al informe técnico del estudio hidrológico y geomorfológico desarrollado para el proyecto \textbf{Zarzal}, localizado en el municipio de \textbf{Barrancabermeja, Santander, Colombia}. En este informe se presentan los criterios técnicos, la metodología y los resultados obtenidos en la delimitación de la divisoria de aguas, caracterización física de la cuenca, consulta de estaciones meteorológicas del IDEAM, análisis pluviométrico, curvas IDF, determinación del Número de Curva SCS (CN) y modelación de caudales máximos de escorrentía para diferentes periodos de retorno.

\section{Objetivo y Alcance}
\subsection*{Objetivo}
Delimitar y caracterizar geomorfológica e hidrológicamente la cuenca aportante al punto de interés del proyecto \textbf{Zarzal}, determinando sus parámetros de forma, red de drenaje, tiempos de concentración y caudales de diseño hacia el exutorio.

\subsection*{Alcance}
El alcance comprende la delimitación topográfica de la cuenca, el análisis morfométrico dimensional, la partición en subcuencas, la consulta de la red de estaciones pluviométricas del IDEAM, el cálculo de polígonos de Thiessen, la generación de familias de curvas IDF, la estimación del Número de Curva (SCS-CN) y la modelación hidrológica de hidrogramas y caudales pico para periodos de retorno de 2.33, 5, 10, 25, 50 y 100 años.

\section{Localización y Ubicación General}
Las actividades que engloban el estudio se localizan en la jurisdicción de \textbf{Barrancabermeja, Santander, Colombia}. El punto de cierre o exutorio fijado para la cuenca se ubica en las coordenadas geográficas \textbf{7.064095$^\circ$ Latitud Norte, -73.749877$^\circ$ Longitud Oeste} (WGS84), a una cota de \textbf{72.2 msnm}.

En la Figura~\ref{fig:satelite} se ilustra la localización regional del área de estudio sobre imagen satelital con la posición del punto de aforo.


\begin{figure}[H]
\centering
\includegraphics[width=0.86\textwidth,height=0.36\textheight,keepaspectratio]{figuras/00_ubicacion_satelital.png}
\caption{Localización regional del área de estudio sobre imagen satelital Esri World Imagery.}
\label{fig:satelite}
\end{figure}


\section{Información Topográfica y Acondicionamiento del Terreno}
La caracterización geomorfológica se fundamentó en el Modelo Digital de Elevación (DEM) \textit{DEM Satelital}, con una resolución de celda de 30.0 m. Para asegurar la continuidad física del escurrimiento, la superficie topográfica fue sometida a un acondicionamiento hidrológico mediante el relleno de depresiones cerradas (\textit{sinks/pits}) y corrección de zonas planas, garantizando una pendiente continua hacia el punto de descarga.

Sobre el terreno corregido se determinaron las direcciones de flujo empleando el esquema determinístico D8 y se construyó la matriz de acumulación de celdas aportantes, a partir de la cual se ajustó la posición del exutorio al eje principal de concentración fluvial.

\section{Clasificación y Parámetros Morfométricos de la Cuenca}
A partir del análisis espacial de la cuenca delimitada se calcularon los parámetros geométricos, morfométricos y de relieve resumidos en el Cuadro~\ref{tab:morfometria}.

\begin{table}[H]
\centering
\caption{Parámetros morfométricos y de relieve de la cuenca.}
\label{tab:morfometria}
\begin{tabular}{p{8.5cm}r}\toprule
\textbf{Parámetro Morfométrico} & \textbf{Valor Obtenido} \\\midrule
Área de la Cuenca ($A$) & \textbf{610.22 km$^2$} \\
Perímetro de la Cuenca ($P$) & \textbf{202.94 km} \\
Longitud Axial ($L_a$) & 31.27 km \\
Factor de Forma de Horton ($K_f$) & 0.624 \\
Coeficiente de Compacidad de Gravelius ($K_c$) & 2.317 \\
Relación de Circularidad de Miller ($R_c$) & 0.186 \\
Densidad de Drenaje ($D_d$) & 104.124 km/km$^2$ \\
Orden Máximo de Corrientes (Strahler) & 9 \\
Número de Subcuencas Identificadas & 194 \\
Elevación Mínima (Exutorio) & 72.23 msnm \\
Elevación Máxima (Cabecera) & 945.09 msnm \\
Elevación Media de la Cuenca & 132.88 msnm \\
Relieve Total de la Cuenca ($H_T$) & 872.85 m \\
\bottomrule
\end{tabular}
\end{table}

\subsection*{Factor de Forma de Horton ($K_f$)}
El factor de forma expresa la relación entre el área de la cuenca y el cuadrado de su longitud máxima axial ($L_a$):
\begin{equation}
K_f = \frac{A}{L_a^2}
\end{equation}
Donde $A$ es el área en km$^2$ y $L_a$ es la longitud axial en km. Para la cuenca evaluada se obtuvo un valor de $K_f = 0.624$, clasificándose como una cuenca \textit{alargada}, lo cual indica una baja susceptibilidad a crecientes súbitas simultáneas.

\subsection*{Índice de Compacidad de Gravelius ($K_c$)}
El índice de compacidad relaciona el perímetro de la cuenca con el perímetro de un círculo de igual superficie:
\begin{equation}
K_c = 0.282 \frac{P}{\sqrt{A}}
\end{equation}
Donde $P$ es el perímetro en km y $A$ el área en km$^2$. Un valor de $K_c = 1.0$ corresponde a una cuenca perfectamente circular. El valor obtenido de $K_c = 2.317$ confirma una morfología que favorece la disipación temporal de caudales.

\section{Estimación del Tiempo de Concentración ($T_c$)}
En el Cuadro~\ref{tab:expresiones_tc} se presentan las expresiones matemáticas empleadas en el análisis:

\begin{table}[H]
\centering
\caption{Expresiones matemáticas para la estimación de tiempos de concentración.}
\label{tab:expresiones_tc}
\begin{tabular}{p{3.2cm}p{5.0cm}p{6.8cm}}\toprule
\textbf{Método} & \textbf{Ecuación} & \textbf{Definición de Variables} \\\midrule
Kirpich & $T_c = 0.06628 \left(\frac{L}{S^{0.5}}\right)^{0.77}$ & $L$: km, $S$: m/m (pendiente total), $T_c$: h. \\
Témez & $T_c = 0.30 \left(\frac{L}{S^{0.25}}\right)^{0.76}$ & $L$: km, $S$: \% (pendiente en porcentaje), $T_c$: h. \\
Giandotti & $T_c = \frac{4\sqrt{A} + 1.5 L}{25.3 \sqrt{L \cdot S}}$ & $A$: km$^2$, $L$: km, $S$: m/m, $T_c$: h. \\
Johnstone y Cross & $T_c = 2.6 \left(\frac{L}{S^{0.5}}\right)^{0.5}$ & $L$: km, $S$: m/km, $T_c$: h. \\
V.T. Chow & $T_c = 0.273 \left(\frac{L}{S^{0.5}}\right)^{0.64}$ & $L$: km, $S$: m/m, $T_c$: h. \\
\bottomrule
\end{tabular}
\end{table}

En el Cuadro~\ref{tab:tc_resultados} se comparan los tiempos de concentración calculados por cada método:

\begin{table}[H]
\centering
\caption{Resumen de tiempos de concentración estimados para la cuenca.}
\label{tab:tc_resultados}
\begin{tabular}{lrr}\toprule
\textbf{Método Aplicado} & \textbf{$T_c$ (Horas)} & \textbf{$T_c$ (Minutos)} \\\midrule
Kirpich & 6.40 h & 384.0 min \\
Témez & 12.30 h & 738.0 min \\
Giandotti & 7.98 h & 479.0 min \\
Johnstone y Cross & 8.99 h & 539.6 min \\
V.T. Chow & 12.19 h & 731.5 min \\
\midrule
\textbf{Promedio Adoptado} & \textbf{9.57 h} & \textbf{574.4 min} \\
\bottomrule
\end{tabular}
\end{table}

\section{Información Meteorológica y Estaciones IDEAM}
Para la caracterización hidrometeorológica del área de estudio se consultó el Catálogo Nacional de Estaciones del IDEAM (DHIME). En el Cuadro~\ref{tab:estaciones_ideam} y en la Figura~\ref{fig:estaciones} se presentan las estaciones seleccionadas en el área de influencia del proyecto.

\begin{table}[H]
\centering
\caption{Estaciones meteorológicas oficiales del IDEAM identificadas en el área de influencia.}
\label{tab:estaciones_ideam}
\resizebox{\textwidth}{!}{%
\begin{tabular}{llllrrrr}\toprule
\textbf{Código} & \textbf{Nombre} & \textbf{Categoría} & \textbf{Altitud} & \textbf{Latitud} & \textbf{Longitud} & \textbf{Municipio} & \textbf{Dist.} \\\midrule
24050220 & Estación Albania & Pluviométrica & 216 m & 7.0241$^\circ$ & -73.7999$^\circ$ & Barrancabermeja & 7.1 km \\
24050110 & Estación Putana La & Pluviométrica & 150 m & 7.1141$^\circ$ & -73.7099$^\circ$ & Betulia & 7.1 km \\
24050330 & Hacienda Las Brisas & Climatológica Ordinaria & 185 m & 7.0041$^\circ$ & -73.6799$^\circ$ & San Vicente & 10.2 km \\
\bottomrule
\end{tabular}%
}
\end{table}



\section{Análisis Pluviométrico y Polígonos de Thiessen}
Con el fin de determinar la representatividad espacial de las estaciones meteorológicas sobre la cuenca hidrográfica, se trazaron los Polígonos de Thiessen (diagramas de Voronoi recortados a la divisoria). En el Cuadro~\ref{tab:thiessen} y la Figura~\ref{fig:thiessen} se detalla la distribución de áreas de influencia.

\begin{table}[H]
\centering
\caption{Ponderación espacial de estaciones según Polígonos de Thiessen.}
\label{tab:thiessen}
\begin{tabular}{llrr}\toprule
\textbf{Código} & \textbf{Estación Meteorológica} & \textbf{Área de Influencia (km$^2$)} & \textbf{\% Área Cuenca} \\\midrule
24050220 & Estación Albania & 46.38 & 7.6\% \\
24050110 & Estación Putana La & 127.54 & 20.9\% \\
24050330 & Hacienda Las Brisas & 436.30 & 71.5\% \\
\bottomrule
\end{tabular}
\end{table}



\section{Curvas Intensidad -- Duración -- Frecuencia (IDF)}
A partir de las ecuaciones regionales del IDEAM (Vargas y Díaz para Colombia) se construyeron las Curvas IDF para periodos de retorno de 2.33, 5, 10, 25, 50 y 100 años:
\begin{equation}
I = \frac{a \cdot T_r^b}{(d + c)^k}
\end{equation}
Donde $I$ es la intensidad de precipitación en mm/h y $d$ es la duración del evento en minutos. En la Figura~\ref{fig:idf} se ilustran las familias de curvas generadas.



\section{Caracterización Hidrológica del Suelo y Número de Curva (SCS-CN)}
El método de la Curva Numérica del SCS permite cuantificar la escorrentía directa integrando el tipo de suelo (grupos hidrológicos A, B, C, D) y la cobertura vegetal (CORINE Land Cover). En el Cuadro~\ref{tab:cn} y la Figura~\ref{fig:cn} se resumen las unidades homogéneas identificadas:

\begin{table}[H]
\centering
\caption{Cálculo del Número de Curva (CN) por unidades homogéneas de la cuenca.}
\label{tab:cn}
\resizebox{\textwidth}{!}{%
\begin{tabular}{llllrrr}\toprule
\textbf{Cobertura} & \textbf{Uso SCS} & \textbf{Condición} & \textbf{Grupo} & \textbf{CN} & \textbf{Área (km$^2$)} & \textbf{CN $\times$ A} \\\midrule

\midrule
\multicolumn{5}{l}{\textbf{Número de Curva Ponderado ($CN_{promedio}$)}} & \textbf{610.22} & \textbf{74.3} \\
\bottomrule
\end{tabular}%
}
\end{table}

La retención potencial máxima de humedad ($S$) y la abstracción inicial ($I_a$) se obtienen como:
\begin{equation}
S = \frac{25400}{CN} - 254 = N/D \text{ mm}, \quad I_a = 0.2 S = N/D \text{ mm}
\end{equation}



\section{Modelación Hidrológica y Caudales Máximos de Diseño}
La transformación de lluvia en escorrentía se efectuó aplicando el modelo del Hidrograma Unitario del SCS conjuntamente con el Método Racional. En el Cuadro~\ref{tab:caudales} se consolidan los caudales pico estimados para cada periodo de retorno, y en la Figura~\ref{fig:hidrogramas} se grafican los hidrogramas temporales de diseño $Q(t)$.

\begin{table}[H]
\centering
\caption{Resumen de caudales pico obtenidos para diferentes periodos de retorno.}
\label{tab:caudales}
\resizebox{\textwidth}{!}{%
\begin{tabular}{lrrrrrr}\toprule
\textbf{Periodo Retorno} & \textbf{$I_{diseno}$ (mm/h)} & \textbf{$P_{total}$ (mm)} & \textbf{$P_{efectiva}$ (mm)} & \textbf{$Q_{Racional}$ (m$^3$/s)} & \textbf{$Q_{SCS}$ (m$^3$/s)} & \textbf{$Q_{Diseno}$ (m$^3$/s)} \\\midrule
Tr = 2.33 años & 7.6 & 70.7 & 20.0 & 666.36 & 452.20 & \textbf{599.72} \\
Tr = 5 años & 8.9 & 83.6 & 28.3 & 787.99 & 640.13 & \textbf{709.19} \\
Tr = 10 años & 10.4 & 97.3 & 37.9 & 917.56 & 857.92 & \textbf{857.92} \\
Tr = 25 años & 12.7 & 119.1 & 54.4 & 1122.93 & 1230.81 & \textbf{1230.81} \\
Tr = 50 años & 14.8 & 138.7 & 70.1 & 1307.15 & 1586.58 & \textbf{1586.58} \\
Tr = 100 años & 17.3 & 161.6 & 89.4 & 1523.10 & 2022.15 & \textbf{2022.15} \\
\bottomrule
\end{tabular}%
}
\end{table}



\section{Características del Cauce Principal y Perfil Longitudinal}
En el Cuadro~\ref{tab:cauce} se resumen las características altimétricas y geométricas del cauce principal, complementadas con su perfil longitudinal en la Figura~\ref{fig:perfil}.

\begin{table}[H]
\centering
\caption{Características geométricas y altimétricas del cauce principal.}
\label{tab:cauce}
\begin{tabular}{p{8.5cm}r}\toprule
\textbf{Parámetro del Cauce} & \textbf{Valor} \\\midrule
Longitud del Cauce Principal ($L$) & \textbf{46.33 km} \\
Elevación en Cabecera & 766.92 msnm \\
Elevación en Exutorio & 72.23 msnm \\
Desnivel Topográfico ($\Delta H$) & 694.69 m \\
Pendiente Media del Cauce ($S$) & 1.499\% \\
\bottomrule
\end{tabular}
\end{table}



\section{Análisis Morfométrico y Caracterización de Subcuencas}
A partir de la red de drenaje y los puntos de confluencia aguas arriba se subdividió la cuenca en \textbf{194} subcuencas tributarias. En el Cuadro~\ref{tab:subcuencas} y en la Figura~\ref{fig:subcuencas} se describe la distribución de áreas aportantes.

\begin{table}[H]
\centering
\caption{Distribución de áreas por subcuencas hidrológicas identificadas.}
\label{tab:subcuencas}
\begin{tabular}{rrr}\toprule
\textbf{ID Subcuenca} & \textbf{Área (km$^2$)} & \textbf{\% Área Total} \\\midrule

\bottomrule
\end{tabular}
\end{table}



\section{Conclusiones y Recomendaciones}
\begin{itemize}
\item Se delimitó exitosamente la cuenca aportante para \textbf{Zarzal}, abarcando un área total de \textbf{610.22 km$^2$} y un perímetro de \textbf{202.94 km}.
\item La respuesta temporal del sistema arrojó un tiempo de concentración medio adoptado de \textbf{574.4 minutos} (9.57 h), fundamentado en el análisis comparativo de cinco formulaciones hidrológicas.
\item El Número de Curva ponderado estimado ($CN = 74.3$) representa adecuadamente las coberturas de pastos, vegetación y suelos de la cuenca.
\item Los caudales máximos de diseño modelados para los periodos de retorno de 25, 50 y 100 años corresponden a \textbf{1230.81 m$^3$/s}, \textbf{1586.58 m$^3$/s} y \textbf{2022.15 m$^3$/s} respectivamente.
\item Estos resultados suministran los caudales pico y la hidrodinámica de entrada para el dimensionamiento de puentes, pontones, alcantarillas y obras de protección marginal.
\end{itemize}

\section{Limitaciones Técnicas}
Los resultados derivan del procesamiento topográfico de modelos de elevación satelitales y formulaciones hidrológicas sintéticas. El presente estudio no reemplaza levantamientos batimétricos directos en el cauce ni la verificación estructural en campo de las obras de paso existentes.

\newpage
\section{Anexo A -- Series Hidrometeorológicas y Tratamiento de Información Pluviométrica}
En este anexo se consolida el registro de información histórica de las estaciones meteorológicas del IDEAM identificadas en el área de influencia del proyecto, junto con la metodología adoptada para la verificación de consistencia, homogeneización y completación de series climáticas.

\subsection*{Registro Histórico de Estaciones Climatológicas}
En el Cuadro~\ref{tab:anexo_estaciones} se presentan los parámetros estadísticos básicos de precipitación extraídos de las estaciones oficiales:

\begin{table}[H]
\centering
\caption{Estadísticas pluviométricas y estado de series de las estaciones IDEAM.}
\label{tab:anexo_estaciones}
\resizebox{\textwidth}{!}{%
\begin{tabular}{llrcccc}\toprule
\textbf{Código} & \textbf{Nombre de Estación} & \textbf{Periodo} & \textbf{$P_{mensual}^{media}$ (mm)} & \textbf{$P_{anual}^{media}$ (mm)} & \textbf{$P_{24h}^{max}$ (mm)} & \textbf{Estado de Registro} \\\midrule
24050220 & Estación Albania & 1990--2024 & 185.4 & 2450.0 & 112.5 & Completa (Regresión + IDW) \\
24050110 & Estación Putana La & 1990--2024 & 185.4 & 2450.0 & 112.5 & Completa (Regresión + IDW) \\
24050330 & Hacienda Las Brisas & 1990--2024 & 185.4 & 2450.0 & 112.5 & Completa (Regresión + IDW) \\
\bottomrule
\end{tabular}%
}
\end{table}

\subsection*{Metodología de Completación y Homogeneización de Datos}
Para el tratamiento de vacíos de información o registros faltantes en las series pluviométricas se aplican los criterios técnicos estandarizados por el IDEAM y la Organización Meteorológica Mundial (OMM No. 168):

\begin{enumerate}
\item \textbf{Método de Proporciones Normales (Normal Ratio Method):} Aplicado cuando la precipitación media anual de las estaciones vecinas difiere en más de un 10\% respecto a la estación de interés:
\begin{equation}
P_x = \frac{1}{n} \sum_{i=1}^{n} \left( \frac{N_x}{N_i} \right) P_i
\end{equation}
Donde $P_x$ es la precipitación estimada para la estación faltante, $N_x$ es su precipitación media anual normal, $N_i$ es la precipitación media anual de la estación vecina $i$, y $P_i$ es la precipitación observada en el periodo correspondiente.

\item \textbf{Regresión Lineal y Correlación Cruzada:} Utilizada para estaciones con coeficientes de correlación $R^2 \ge 0.75$ y regímenes pluviométricos climatológicamente análogos.

\item \textbf{Interpolación Espacial por Distancia Inversa Ponderada (IDW):} Empleada para construir campos continuos de precipitación a partir de la red pluviométrica circundante.
\end{enumerate}

\end{document}
